%-------------------------------------------------------------------------------
% 2. 2025 MixUP AI Datathon — 국내 기업 철강 공정 데이터 분석 (본선 5위)
%-------------------------------------------------------------------------------
\cvsection{2. 2025 MixUP AI Datathon : 국내 기업 철강 공정 데이터 분석}

\vspace{0.6mm}\noindent{\bfseries\color{awesome}본선 진출 및 최종 5위}\par
\projsummary{스마트팩토리 센서 데이터를 활용한 품질 불량(OK/NG) 조기 탐지 및 공정 최적화 전략 수립}
\projfig{p2.png}

\cvsubsection{배경 및 목표}
\begin{cvitems}
  \item {\textbf{배경}\enskip 철강 제조는 고온·고압 환경에서 수많은 센서 로그가 발생하는 복잡한 공정으로, 미세한 변수 제어 실패가 막대한 자재 손실과 품질 저하로 직결.}
  \item {\textbf{목표}\enskip 스마트팩토리 솔루션 기업 위즈코어(Wizcore)의 실제 철강 공정 데이터를 활용해 양품(OK)/불량품(NG)을 판별하는 이진 분류 모델을 개발하고, 데이터 분석을 통한 공정 최적화 및 운용 전략을 제안.}
  \item {\textbf{개발 기간}\enskip 2025.05.12 -- 2025.05.18 (대면 기간 2025.05.17--18)}
  \item {\textbf{기술 스택}\enskip Python · LightGBM · Optuna · SHAP · RFECV · Scikit-learn}
\end{cvitems}

\cvsubsection{데이터 소개}
\begin{cvitems}
  \item {\textbf{학습/평가 데이터}\enskip 철강 생산 공정의 제품별 센서 로그(장력·스피드 등), 생산지시 정보, 품질 결과(OK/NG).}
  \item {\textbf{메타 데이터}\enskip 공정별 설비 부하 기준(암페어 표준화) 및 제품 규격별 표준 운영 가이드라인.}
\end{cvitems}

\cvsubsection{설계 과정 및 수행 내용}
\begin{cvitems}
  \item {\textbf{① 물리적 인과관계 분석 및 데이터 정합성 검증}\enskip 노이즈·이상치가 잦은 공정 데이터가 실제 물리 현상을 올바르게 반영하는지 모델링 전 검증. '속도가 증가하면 점성 저항이 커져 장력이 선형적으로 증가'한다는 물리 원리로 분석해 두 변수 간 높은 상관($R^2=0.7914$)을 확인.}
  \item {\textbf{② RFECV 기반 변수 최적화 및 도메인 기반 피처 설계}\enskip 27개 다차원 변수 중 노이즈·중복을 제거해 복잡도를 낮추고, RFECV(재귀적 특징 제거)로 기여도 높은 \textbf{최적 6개 변수}(기준장력·장력1·스피드1 등)를 선별. 규격에 구애받지 않는 '공정 상태 지표'로 \textbf{장력비}(기준 대비 편차 비율 = 상대적 공정 부하), \textbf{스피드 표준편차}(롤러 간 속도 불균형 = 설비 동기화 오류), \textbf{측량완전결측 플래그}(장력·스피드 동시 0값 구간의 양품률이 0\%임을 확인해 결측이 아닌 '불량 신호'로 정의)를 설계.}
  \item {\textbf{③ 고정밀 품질 예측 및 데이터 불균형 최적화}\enskip 불량(NG)이 현저히 적은 제조 현장의 불균형으로 생기는 양품 편향을 해결하기 위해 Class Weight(OK:1 / NG:3)를 적용하고, 비선형 패턴 학습에 우수한 LightGBM을 선정. F1-Score와 ROC-AUC를 가중 결합한 대회 전용 지표로 Optuna 하이퍼파라미터 튜닝.}
  \item {\textbf{④ XAI(SHAP)를 이용한 공정 최적화 전략 수립}\enskip 불량을 맞히는 데 그치지 않고 현장 관리자가 즉각 조치할 수 있도록 판단 근거를 시각화. SHAP 분석으로 '스피드가 빠를수록 기준 장력 설정의 미세한 오차가 품질에 훨씬 민감하게 작용'하는 변수 간 교호작용을 도출하고, 이를 '고위험 구간 자동 경고 시스템'과 '실시간 피드백 제어 로직' 도입으로 비즈니스 전략화.}
  \item {\textbf{⑤ 모델 학습 시 '타겟 누수(Target Leakage)' 방지}\enskip 일자·생산지시번호 등 예측 시점에는 알 수 없거나 정답과 직결된 식별자·비즈니스 변수 14개 컬럼을 사전에 식별·제거해, 실제 현장 투입 시에도 동일한 성능을 유지하는 데이터 파이프라인을 구축.}
\end{cvitems}

\cvsubsection{성과}
\begin{cvparagraph}
본선 진출 및 \textbf{최종 5위}. 예측을 넘어 SHAP 기반 공정 최적화 전략까지 제안해 실무 적용 가능성을 인정받음.
\end{cvparagraph}
